% Source printed pp. 431-437 / PDF pp. 439-445.
\section{Kempf 最优不稳定化定理}\label{src:8.6}
\begin{epigraph}
事实上, Tits 猜想旗复形的每个凸子集 $C$ 都有自然的中心 $\delta$. 此时 $P(\delta)$ 就是``最坏的''抛物子群.
\epigraphauthor{David Mumford \cite[第 64 页]{GIT}}
\end{epigraph}
由 Hilbert--Mumford 判据 (\bookref{8.4.4}), 对每个非半稳定点 $x\in X$, 存在单参数子群 $\lambda$, 使 Hilbert--Mumford 指数 $\mu(x,\lambda)<0$. Kempf 最优不稳定化定理 (\bookref{8.6.6}) 断言存在使归一化 Hilbert--Mumford 指数 $\mu(x,\lambda)/\|\lambda\|$ 最小的典范单参数子群 $\lambda$. 换言之, 存在对 $x$ 的不稳定性负有最大责任的``最坏''单参数子群.
\subsection{衡量不稳定性}\label{src:8.6.1-subsection}
给定代数闭域 $\Bbbk$ 上的代数群 $G$, 以 $\mathbb X_*(G)$ 表示单参数子群 $\mathbb G_m\to G$ 的集合. 回忆对环面 $T\cong\mathbb G_m^n$, 有 $\mathbb X_*(T)\cong\mathbb Z^n$ (见例~\bookref{B.1.25}).
\begin{definition}\label{src:8.6.1}
$\mathbb X_*(G)$ 上的\newterm{长度}{length} $\|-\|$ 是其上的非负实值函数, 满足共轭不变性, 即对 $\lambda\in\mathbb X_*(G)$ 与 $g\in G(\Bbbk)$ 有 $\|g\lambda g^{-1}\|=\|\lambda\|$; 且对每个极大环面 $T\subseteq G$, $\mathbb X_*(T)$ 上存在正定整值双线性形式 $(-,-)$, 使对 $\lambda\in\mathbb X_*(T)$ 有 $(\lambda,\lambda)=\|\lambda\|^2$.
\end{definition}
\begin{example}\label{src:8.6.2}
若 $G=\mathrm{GL}_n$, 则任意单参数子群 $\lambda$ 共轭于形如 $t\mapsto\operatorname{diag}(t^{d_1},\ldots,t^{d_n})$ 的单参数子群, 可定义 $\|\lambda\|=\sqrt{d_1^2+\cdots+d_n^2}$.
\end{example}
\begin{example}\label{src:8.6.3}
对每个约化群 $G$, $\mathbb X_*(G)$ 上都存在长度 $\|-\|$. 为证明此事, 取极大环面 $T\subseteq G$, 在 $\mathbb X_*(T)$ 上选择关于 Weyl 群 $W:=N(T)/T$ 的共轭作用不变的正定整值双线性形式 $(-,-)$. $\mathbb X_*(G)$ 在 $G$ 下的共轭类与 $\mathbb X_*(T)$ 在 $W$ 下的共轭类之间有双射 $\mathbb X_*(G)/G\cong\mathbb X_*(T)/W$. 换言之, 对每个 $\lambda\in\mathbb X_*(G)$, 存在 $g\in G(\Bbbk)$, 使 $g\lambda g^{-1}\in\mathbb X_*(T)$; 对任何也满足 $g'\lambda g'^{-1}\in\mathbb X_*(T)$ 的 $g'\in G(\Bbbk)$, $g\lambda g^{-1}$ 与 $g'\lambda g'^{-1}$ 在 $W$ 下共轭. 所以 $\|\lambda\|^2:=(g\lambda g^{-1},g\lambda g^{-1})$ 良定义.
\end{example}
设 $X=\Spec A$ 是带 $G$-作用的仿射 $\Bbbk$-概形, $x_0\in X(\Bbbk)$ 是轨道闭的点. 对满足 $Gx_0\subseteq\overline{Gx}$ 的点 $x\in X(\Bbbk)$ 和使 $\lim_{t\to0}\lambda(t)\cdot x$ 存在的单参数子群 $\lambda:\mathbb G_m\to G$, 定义 $x$ 关于 $\lambda$ 的 Hilbert--Mumford 指数为
\begin{equation}\label{src:8.6.4}
\mu(x,\lambda)=-\deg f_{x,\lambda}^{-1}(Gx_0),
\end{equation}
其中 $f_{x,\lambda}:\mathbb A^1\to X$ 是映射 $\mathbb G_m\to X$, $t\mapsto\lambda(t)\cdot x$ 的延拓. 注意若 $\lim_{t\to0}\lambda(t)\cdot x\notin Gx_0$, 则 $\mu(x,\lambda)=0$.

在 $G$-等变嵌入 $X\hookrightarrow\mathbb P(V)$ 的射影情形中, 我们已在 \bookeqref{8.4.2} 中定义 Hilbert--Mumford 指数 $\mu(x,\lambda)$: 选择 $V$ 的基, 使 $\mathbb G_m$ 在 $\mathbb A(V)=\mathbb A^n$ 上以权 $d_1,\ldots,d_n$ 作用, 并选择 $x$ 的提升 $\widehat x=(u_1,\ldots,u_n)\in\mathbb A(V)$, 则 $-\mu(x,\lambda)$ 是满足 $u_i\ne0$ 的最小 $d_i$. 若 $\lim_{t\to0}\lambda(t)\cdot\widehat x$ 存在, 则这与 \bookeqref{8.6.4} 的定义一致. 为看出此事, 注意 $t\mapsto\lambda(t)\cdot\widehat x$ 的延拓 $f_{\widehat x,\lambda}:\mathbb A^1\to\mathbb A^n$ 是映射 $t\mapsto(t^{d_i}u_i)$, 且 $f_{\widehat x,\lambda}^{-1}(0)=\Spec\Bbbk[t]/(t^d)$, 其中 $d$ 为满足 $u_i\ne0$ 的最小 $d_i$.

因为 $\mu(x,\lambda^n)=n\mu(x,\lambda)$, 自然考虑\newterm{归一化 Hilbert--Mumford 指数}{normalized Hilbert--Mumford index}
\[
\frac{\mu(x,\lambda)}{\|\lambda\|},
\]
以衡量 $\lambda(t)\cdot x$ 趋近闭轨道 $Gx_0$ 的速度. 此指数越负, $\lambda(t)\cdot x$ 趋近 $Gx_0$ 越快. Kempf 证明存在使此指数最小的单参数子群, 且它在共轭意义下唯一.
\subsection{Kempf 定理的陈述}\label{src:8.6.2-subsection}
给定单参数子群 $\lambda:\mathbb G_m\to G$, 回忆 \S\bookref{B.1.4-subsection} 中, 中心化子 $C_\lambda$ 是与 $\lambda$ 交换的元素所成的子群, 抛物子群为 $P_\lambda=\{g\in G\mid\lim_{t\to0}\lambda_0(t)g\lambda_0(t)^{-1}\text{ 存在}\}$, 幺幂子群为 $U_{\lambda_0}=\{g\in G\mid\lim_{t\to0}\lambda_0(t)g\lambda_0(t)^{-1}=1\}$. 有短正合列 $1\to U_\lambda\to P_\lambda\to C_\lambda\to1$, 将 $U_\lambda$ 识别为 $P_\lambda$ 的幺幂根.
\begin{theorem}[Kempf 最优不稳定化定理, 仿射版本]\label{src:8.6.5}
设 $G$ 是代数闭域 $\Bbbk$ 上的约化群, $\mathbb X_*(G)$ 配有长度 $\|-\|$. 设 $X=\Spec A$ 是带 $G$-作用的 $\Bbbk$ 上有限型仿射概形, $x_0\in X(\Bbbk)$ 的轨道闭. 对每个满足 $Gx_0\subseteq\overline{Gx}$ 的点 $x\in X(\Bbbk)$, 存在单参数子群 $\lambda_0:\mathbb G_m\to G$, 使 $\mu(x,\lambda_0)/\|\lambda_0\|$ 达到最小值
\[
M(x):=\inf_{\lambda\in\mathbb X_*(G)}\left\{\frac{\mu(x,\lambda)}{\|\lambda\|}\ \middle|\ \lim_{t\to0}\lambda(t)\cdot x\in Gx_0\right\}.
\]
若 $\lambda_0'$ 是另一个这样的单参数子群, 则 $P_{\lambda_0}=P_{\lambda_0'}$, 且存在唯一元素 $u\in U_{\lambda_0}$, 使 $\lambda_0'=u\lambda_0u^{-1}$. 每个极大环面 $T\subseteq P_{\lambda_0}$ 都含有唯一达到此最小值的不可分元素.
\end{theorem}
将仿射版本 (定理~\bookref{8.6.5}) 应用于非半稳定点 $x\in\mathbb P(V)$ 的提升 $\widehat x\in\mathbb A(V)$, 即得到下面的射影版本; 此时 $\overline{G\widehat x}$ 中的闭轨道是不动点 $0$. 射影版本对约化群也成立, 但因为我们只在线性约化群的情形下讨论了半稳定性, 这里限制在线性约化群的情形.
\begin{theorem}[Kempf 最优不稳定化定理, 射影版本]\label{src:8.6.6}
设 $G$ 是代数闭域 $\Bbbk$ 上的约化群, $\mathbb X_*(G)$ 配有长度 $\|-\|$. 设 $X\subseteq\mathbb P(V)$ 是 $G$-等变闭子概形, 其中 $V$ 是有限维 $G$-表示. 对每个非半稳定点 $x\in X(\Bbbk)$, 存在单参数子群 $\lambda_0:\mathbb G_m\to G$, 使 $\mu(x,\lambda_0)/\|\lambda_0\|$ 达到最小值
\[
M(x):=\inf_{\lambda\in\mathbb X_*(G)}\left\{\frac{\mu(x,\lambda)}{\|\lambda\|}\right\}.
\]
若 $\lambda_0'$ 是另一个这样的单参数子群, 则 $P_{\lambda_0}=P_{\lambda_0'}$, 且存在唯一元素 $u\in U_{\lambda_0}$, 使 $\lambda_0'=u\lambda_0u^{-1}$. 每个极大环面 $T\subseteq P_{\lambda_0}$ 都含有唯一达到此最小值的不可分元素.
\end{theorem}
\begin{definition}\label{src:8.6.7}
满足定理~\bookref{8.6.5} 或定理~\bookref{8.6.6} 的任意 $\lambda_0$ 称为 $x$ 的\newterm{最优不稳定化单参数子群}{optimal destabilizing one-parameter subgroup}, $M(x)<0$ 称为 $x$ 的\newterm{最优归一化 Hilbert--Mumford 指数}{optimal normalized Hilbert--Mumford index}.
\end{definition}
\subsection{Kempf 定理的证明}\label{src:8.6.3-subsection}
当 $x_0\in\overline{Gx}$ 是不动点时, 下面的证明较简单, 例如射影版本中闭轨道为 $0\in\mathbb A(V)$ 的情形. 读者不妨始终将此情形记在心中.
\begin{proof}[定理~\bookref{8.6.5} 的证明]
由命题~\bookref{B.1.21}, 可选取有限维 $G$-表示 $V$, $W$ 和 $G$-等变映射
\begin{equation}\label{src:8.6.8}
\begin{tikzcd}[ampersand replacement=\&]
{}\&\mathbb A(V)\\
X\arrow[ur,hook,"i"]\arrow[dr,"f"']\&{}\\
{}\&\mathbb A(W),
\end{tikzcd}
\end{equation}
其中 $i:X\hookrightarrow\mathbb A(V)$ 是满足 $i(x_0)=0$ 的闭浸入, $f:X\to\mathbb A(W)$ 满足 $f^{-1}(0)=Gx_0$. (当 $x_0$ 为不动点时, 可在 \bookeqref{8.6.8} 中取 $f=i$.) 单参数子群 $\lambda:\mathbb G_m\to G$ 在 $V$, $W$ 上诱导 $\mathbb G_m$-作用, 因而有分次 $V=\bigoplus_{d\in\mathbb Z}V_d$, $W=\bigoplus_{d\in\mathbb Z}W_d$. 定义
\begin{align*}
m(i(x),\lambda)&=\min\{d\mid i(x)\text{ 在 }V_d\text{ 上的投影非零}\},\\
m(f(x),\lambda)&=\min\{d\mid f(x)\text{ 在 }W_d\text{ 上的投影非零}\}.
\end{align*}
对任意 $g\in G$, 有等式 $m(i(x),\lambda)=m(i(g\cdot x),g\lambda g^{-1})$ 与 $m(f(x),\lambda)=m(f(g\cdot x),g\lambda g^{-1})$. 容易看出, 若 $\lim_{t\to0}\lambda(t)\cdot x$ 存在, 则 $\mu(x,\lambda)=-m(f(x),\lambda)$, 且
\begin{align*}
\lim_{t\to0}\lambda(t)\cdot x\text{ 存在}&\quad\Longleftrightarrow\quad m(i(x),\lambda)\geq0,\\
\lim_{t\to0}\lambda(t)\cdot x\in Gx_0&\quad\Longleftrightarrow\quad m(i(x),\lambda)\geq0\text{ 且 }m(f(x),\lambda)>0.
\end{align*}
由不稳定化定理 (\bookref{8.3.6}), 存在 $\lambda_x\in\mathbb X_*(G)$, 使 $m(i(x),\lambda_x)\geq0$ 且 $m(f(x),\lambda_x)>0$.

\textit{环面的情形:} 设 $T\subseteq G$ 是包含 $\lambda_x$ 的极大环面. 可将 $T$-表示分解为 $V=\bigoplus_{\chi\in\mathbb X^*(T)}V_\chi$, 其中 $\mathbb X^*(T)$ 为 $T$ 的特征标集. $i(x)\in V$ 关于 $T$ 的\newterm{状态}{state}是集合
\[
\operatorname{State}_T(i(x)):=\{\chi\in\mathbb X^*(T)\mid i(x)\text{ 在 }V_\chi\text{ 上的投影非零}\}.
\]
类似地, 有 $f(x)\in W$ 关于 $T$ 的状态 $\operatorname{State}_T(f(x))\subseteq\mathbb X^*(T)$. 以 $\langle-,-\rangle$ 表示自然配对 $\mathbb X^*(T)\times\mathbb X_*(T)\to\mathbb Z$. 对单参数子群 $\lambda\in\mathbb X_*(T)$, 有
\begin{align*}
m(i(x),\lambda)&=\min_{\chi\in\operatorname{State}_T(i(x))}\langle\chi,\lambda\rangle,\\
m(f(x),\lambda)&=\min_{\chi\in\operatorname{State}_T(f(x))}\langle\chi,\lambda\rangle.
\end{align*}
我们断言函数 $\lambda\mapsto m(f(x),\lambda)/\|\lambda\|$ 在集合 $\{\lambda\ne0\in\mathbb X_*(T)\mid m(i(x),\lambda_T)\geq0\}$ 上于某个单参数子群 $\lambda_T$ 达到最大值, 且任何其他达到此最小值的单参数子群都是 $\lambda_T$ 的正倍数. 将下面的引理~\bookref{8.6.9} 应用于格 $L=\mathbb X_*(T)\cong\mathbb Z^r$ 以及 $\mathbb X^*(T)\cong\Hom_{\mathbb Z}(L,\mathbb Z)$ 的子集 $F:=\operatorname{State}_T(i(x))$, $G:=\operatorname{State}_T(f(x))$, 得到的恰是此结论.

\textit{一般情形:} 对每个极大环面 $T\subseteq G$ 与每个 $g\in G(\Bbbk)$, 有由 $\chi\leftrightarrow g\chi g^{-1}$ 给出的等同 $\mathbb X^*(T)\cong\mathbb X^*(gTg^{-1})$, 在此等同下 $\operatorname{State}_T(i(x))=\operatorname{State}_{gTg^{-1}}(i(gx))$. 对单参数子群 $\lambda\in\mathbb X_*(G)$, 我们已看到对 $g\in G(\Bbbk)$ 有 $m(f(x),\lambda)=m(f(gx),g\lambda g^{-1})$. 事实上, 我们断言对 $p\in P_\lambda$ 有 $m(f(x),\lambda)=m(f(x),p\lambda p^{-1})$. 由对称性, 只须证明 $m(f(x),\lambda)\leq m(f(x),p\lambda p^{-1})$. 将 $-m(f(x),\lambda)$ 理解为使 $\lim_{t\to0}t^d\lambda(t)\cdot f(x)\in\mathbb A(W)$ 存在的最小整数 $d$, 只须证明 $\lim_{t\to0}t^dp\lambda(t)p^{-1}\cdot f(x)\in\mathbb A(W)$ 存在. 这由计算
\begin{align*}
\lim_{t\to0}\bigl(t^dp\lambda(t)p^{-1}\cdot f(x)\bigr)
&=\lim_{t\to0}\Bigl(p\cdot(\lambda(t)p^{-1}\lambda(t)^{-1})\cdot(t^d\lambda(t)f(x))\Bigr)\\
&=p\cdot\Bigl(\lim_{t\to0}\lambda(t)p^{-1}\lambda(t)^{-1}\Bigr)\cdot\Bigl(\lim_{t\to0}t^d\lambda(t)f(x)\Bigr)
\end{align*}
得到.

现在证明函数 $\lambda\mapsto m(f(x),\lambda)/\|\lambda\|$ 在
\[
\Lambda:=\{\lambda\in\mathbb X_*(G)\mid m(i(x),\lambda)\geq0\}
\]
上达到最小值. 由环面的情形, 对每个 $g\in G(\Bbbk)$, 每个非空集 $\mathbb X_*(gTg^{-1})\cap\Lambda$ 上都有最小值, 此最小值由 $\mathbb X_*(T)$ 的子集 $\operatorname{State}_{gTg^{-1}}(i(x))\cong\operatorname{State}_T(i(g^{-1}u))$ 与 $\operatorname{State}_{gTg^{-1}}(f(x))\cong\operatorname{State}_T(f(g^{-1}u))$ 决定. 这些子集包含于满足 $V_\chi\ne0$ (分别地, $W_\chi\ne0$) 的特征标组成的有限集中, 故 $g$ 遍历 $G(\Bbbk)$ 时只有有限多个最小值. 任意 $\lambda\in\mathbb X_*(G)$ 的像均包含于某个 $gTg^{-1}$ 中, 因而存在由某个单参数子群 $\lambda_0\in\Lambda$ 达到的全局最小值. 可设 $\lambda_0$ 不可分, 即不能写成另一单参数子群的正倍数.

为证明唯一性, 取包含 $\lambda_0$ 的极大环面 $T_0\subseteq G$. 由环面的情形, $\lambda_0\in\mathbb X_*(T_0)\cap\Lambda$ 是达到最小值的唯一不可分单参数子群. 对 $p\in P_{\lambda_0}$, 共轭单参数子群 $p\lambda_0p^{-1}$ 也达到最小值. 因为其他任意极大环面 $T'\subseteq P_{\lambda_0}$ 都形如 $pT_0p^{-1}$, $p\in P_{\lambda_0}$, 所以 $\mathbb X_*(T')\cap\Lambda$ 也含有达到最小值的唯一不可分元素. 最后, 设 $\lambda_1\in\mathbb X_*(G)$ 是另一个达到最小值的不可分元素. 交 $P_{\lambda_0}\cap P_{\lambda_1}$ 包含 $G$ 的某个极大环面 $T$ (命题~\bookref{B.1.29}(d)), 可写 $\lambda_T=p_0\lambda_0p_0^{-1}=p_1\lambda_1p_1^{-1}$, $p_0,p_1\in P_{\lambda_T}$. 于是 $P_{\lambda_0}=P_{\lambda_T}=P_{\lambda_1}$, 且 $\lambda_0$ 与 $\lambda_1$ 由 $U_{\lambda_T}$ 的唯一元素共轭 (命题~\bookref{B.1.29}(c)). 另见 \cite[定理 3.4]{Kem78}.
\end{proof}
上述论证使用了下面的凸几何引理.
\begin{lemma}\label{src:8.6.9}
设 $\Lambda$ 是有限维格, $F$, $G$ 是 $\Lambda^\vee=\Hom_{\mathbb Z}(\Lambda,\mathbb Z)$ 的非空有限子集. 假定 $\Lambda_{\mathbb R}=\Lambda\otimes_{\mathbb Z}\mathbb R$ 上有在 $\Lambda$ 上取整值的正定内积. 定义
\begin{align*}
f_{\min}:\Lambda_{\mathbb R}&\longrightarrow\mathbb R,\quad\lambda\longmapsto\min_{f\in F}f(\lambda),\\
g_{\min}:\Lambda_{\mathbb R}&\longrightarrow\mathbb R,\quad\lambda\longmapsto\min_{g\in G}g(\lambda).
\end{align*}
设存在 $\lambda\in\Lambda_{\mathbb R}$ 使 $f_{\min}(\lambda)\geq0$ 且 $g_{\min}(\lambda)>0$. 则函数
\[
\begin{aligned}
C_F:=\{\lambda\ne0\in\Lambda_{\mathbb R}\mid f_{\min}(\lambda)\geq0\}&\longrightarrow\mathbb R,\\
\lambda&\longmapsto g_{\min}(\lambda)/\|\lambda\|
\end{aligned}
\]
达到最大值 $M$. 存在唯一元素 $\lambda_0\in C_F\cap\Lambda$, 使 $M=g_{\min}(\lambda_0)/\|\lambda_0\|$, 且任何满足 $M=g_{\min}(\lambda)/\|\lambda\|$ 的其他元素 $\lambda\in C_F\cap\Lambda$ 都是 $\lambda_0$ 的整数倍.
\end{lemma}
\begin{proof}
集合 $\{\lambda\in C_F\mid g_{\min}(\lambda)\geq1\}$ 闭且凸, 故包含距原点最近的唯一点 $\lambda'$. 因为对 $\alpha\in\mathbb R$ 有 $g_{\min}(\alpha\lambda')=\alpha g_{\min}(\lambda')$, 必有 $g_{\min}(\lambda')=1$, 且 $\lambda'\in C_F$ 是满足 $g_{\min}(\lambda')=1$, $g_{\min}(\lambda')/\|\lambda'\|=M$ 的唯一点.

现在证明 $\lambda'$ 张成的射线包含整点. 若 $\lambda'$ 位于 $\{\lambda\in C_F\mid g_{\min}(\lambda)\geq1\}$ 的内部, 即对所有 $f\in F$ 有 $f(\lambda')>0$, 且有唯一 $g\in G$ 使 $g(\lambda')=1$, 则 $\lambda'$ 是由 $g=1$ 定义的仿射平面上距原点最近的点. 我们断言 $\lambda'=g^*/\langle g^*,g^*\rangle$, 其中 $g^*\in\Lambda_{\mathbb R}$ 是使对所有 $\lambda\in\Lambda_{\mathbb R}$ 有 $\langle g^*,\lambda\rangle=g(\lambda)$ 的唯一点. 事实上, 点 $\lambda'$ 位于平面 $g=1$ 上; 对平面上其他任意点 $\lambda$, 有 $\langle\lambda',\lambda\rangle=1/\langle g^*,g^*\rangle=\langle\lambda',\lambda'\rangle$. Cauchy--Schwarz 不等式给出 $\langle\lambda',\lambda'\rangle^2=\langle\lambda',\lambda\rangle^2\leq\langle\lambda',\lambda'\rangle\langle\lambda,\lambda\rangle$, 所以 $\langle\lambda',\lambda'\rangle\leq\langle\lambda,\lambda\rangle$. 因为内积和 $g$ 均取整值, 有 $g^*\in\Lambda$. 取 $g^*$ 张成的射线上唯一的不可分元素为 $\lambda_0$.

为归约到此情形, 设 $f_1,\ldots,f_t\in F$ 是满足 $f_i(\lambda')=0$ 的函数, $g_1,\ldots,g_s\in G$ 是满足 $g_i(\lambda')=g_{\min}(\lambda')$ 的函数. 因为每个 $f_i$ 和 $g_j$ 均取整值, 可限制到子空间
\[
W:=\left\{\lambda\in\Lambda_{\mathbb R}\ \middle|\ \begin{aligned}f_1(\lambda)&=\cdots=f_t(\lambda)=0,\\g_1(\lambda)&=\cdots=g_s(\lambda)\end{aligned}\right\}
\]
及格 $W\cap\Lambda$. 则 $\lambda'$ 位于 $\{\lambda\in C_F\cap W\mid g_{\min}(\lambda)\geq1\}$ 的内部, 因而是由 $g_1=1$ 定义的仿射平面上距原点最近的点.
\end{proof}
\begin{corollary}\label{src:8.6.10}
在定理~\bookref{8.6.5} 或定理~\bookref{8.6.6} 的设置中, 存在唯一态射 $f:[\mathbb A^1/\mathbb G_m]\to[X/G]$, 满足 $f(1)\simeq x$ 且 $f(0)\simeq x_0$.
\end{corollary}
\begin{proof}
由命题~\bookref{7.9.1}, 态射 $[\mathbb A^1/\mathbb G_m]\to[X/G]$ 由满足 $\lim_{t\to0}\lambda(t)x\in Gx_0$ 的单参数子群 $\lambda$ 决定, $\lambda$ 在 $P_\lambda$ 的共轭意义下唯一. 任意两个最优不稳定化单参数子群在 $U_\lambda$ (因而在 $P_\lambda$) 下共轭, 故由从 $[\mathbb A^1/\mathbb G_m]$ 到商叠的映射的解释 (命题~\bookref{7.9.1}) 得到结论.
\end{proof}

\subsection{例子与推广}\label{src:8.6.4-subsection}
\begin{example}\label{src:8.6.11}
重新考察 $\mathrm{SL}_2$ 在 $(\mathbb P^1)^n$ 上的作用, 线性化由 Segre 嵌入 $(\mathbb P^1)^n\hookrightarrow\mathbb P^{2^n-1}$ 给出 (例~\bookref{8.5.1}). 非半稳定轨迹由超过 $n/2$ 个点相等的元组 $x=(p_1,\ldots,p_n)$ 组成. 假定恰有 $k>n/2$ 个点相等. 因为 Hilbert--Mumford 指数是对称的, 可假定前 $k$ 个点相等. 若 $\lambda:\mathbb G_m\to\mathrm{SL}_2$ 是单参数子群, 可取 $g\in\mathrm{SL}_2(\Bbbk)$ 使 $g\lambda g^{-1}=\lambda_0^d$, 其中 $d\in\mathbb Z$, $\lambda_0(t)=\operatorname{diag}(t^{-1},t)$. 重新缩放范数后可假定 $\|\lambda_0\|=1$. 还假定 $d\geq0$, 因为 $d<0$ 的情形可类似处理. 于是
\[
\frac{\mu(x,\lambda)}{\|\lambda\|}=\frac{\mu(gx,g\lambda g^{-1})}{\|g\lambda g^{-1}\|}=\mu(gx,\lambda_0).
\]
此指数为负当且仅当 $gx=\{0,\ldots,0,p_{k+1},\ldots,p_n\}$, 此时 $\mu(gx,\lambda_0)=n-2k$. 因而 $\lambda_0$ (分别地, $g^{-1}\lambda_0g$) 是 $gx$ (分别地, $x$) 的最优不稳定化单参数子群. 注意, 由下三角矩阵组成的抛物子群 $P_{\lambda_0}\subseteq\mathrm{SL}_2$ 也是 $0\in\mathbb P^1$ 的稳定子, 所以 $G_{gx}\subseteq P_{\lambda_0}$.
\end{example}
\begin{exercise}\label{src:8.6.12}
在 $\mathrm{SL}_3$ 对三元三次型空间 $\mathbb P(\Gamma(\mathbb P^2,\mathcal O(3)))$ 的作用下, 尖点三次曲线 $y^2z-x^3$ 不是半稳定的 (见习题~\bookref{8.5.3}). 求最优不稳定化单参数子群.
\end{exercise}
\begin{exercise}\label{src:8.6.13}
设 $G$ 是代数闭域 $\Bbbk$ 上的约化群, 在 $\mathbb X_*(G)$ 上给定长度 $\|{-}\|$. 设 $X=\Spec A$ 是带有 $G$ 作用的有限型仿射 $\Bbbk$-概形, $x_0\in X(\Bbbk)$ 的 $G$-轨道闭. 设 $x\in X(\Bbbk)$ 满足 $Gx_0\subseteq\overline{Gx}$, $P_x$ 是由 Kempf 最优不稳定化定理 (\bookref{8.6.5}) 确定的抛物子群.
\begin{enumerate}[(a)]
\item 证明对所有 $g\in G(\Bbbk)$ 有 $gP_xg^{-1}=P_{gx}$.

\textit{提示:} 证明若某单参数子群 $\lambda$ 满足 $P_x=P_\lambda$, 则 $P_{gx}=P_{g\lambda g^{-1}}$.
\item 证明 $G_x\subseteq P_x$.

\textit{提示:} 使用抛物子群 $P$ 满足 $N_G(P)=P$ (命题~\bookref{B.1.29}).
\end{enumerate}
\end{exercise}
下面的判据有时可以通过仅计算一个固定极大环面中单参数子群的 Hilbert--Mumford 指数来检验稳定性或半稳定性.
\begin{exercise}[Kempf--Morrison 判据]\label{src:8.6.14}
设 $G=\mathrm{GL}(W)$ 或 $\mathrm{SL}(W)$, 其中 $W$ 是特征 $0$ 的代数闭域 $\Bbbk$ 上的有限维向量空间. 设 $X\subseteq\mathbb P(V)$ 是 $G$-不变闭子概形, $V$ 是有限维 $G$-表示, $x\in X(\Bbbk)$. 假定存在线性约化子群 $H\subseteq G_x$, 使 $W$ 分解为互异 $H$-表示的直和. 设 $T\subseteq G$ 是与此分解相容的极大环面. 证明
\begin{align*}
x\in X^{\mathrm{ss}}&\iff\text{对所有 $\lambda:\mathbb G_m\to T$ 有 }\mu(x,\lambda)\leq0,\\
x\in X^{\mathrm s}&\iff\text{对所有 $\lambda:\mathbb G_m\to T$ 有 }\mu(x,\lambda)<0,
\end{align*}
且若 $x$ 不是半稳定的, 则 $T$ 包含最优不稳定化单参数子群.

\textit{提示:} 若 $x\notin X^{\mathrm{ss}}$, 设 $\lambda_0:\mathbb G_m\to G$ 是最优不稳定化单参数子群, $0=V_0\subseteq V_1\subseteq\cdots\subseteq V_k=V$ 是抛物子群 $P_{\lambda_0}$ 诱导的滤过. 使用习题~\bookref{8.6.13} 推出每个 $V_i$ 都是 $H$-不变的, 并使用关于 $H$-表示 $V$ 的假设证明每个 $V_i$ 都是 $T$-不变的; 因而 $T\subseteq P_{\lambda_0}$. 再次应用 Kempf 最优不稳定化定理, 在 $T$ 中找到满足 $\mu(x,\lambda)<0$ 的 $\lambda$. 若 $x\notin X^{\mathrm s}$, 取 $x$ 的提升 $\widehat x\in\mathbb A(V)$ 及具有闭轨道的点 $\widehat x_0\in\overline{G\widehat x}$, 使用 Kempf 最优不稳定化定理的仿射版本重复上述论证.
\end{exercise}
\begin{exercise}[完美域上不稳定化单参数子群的存在性]\label{src:8.6.15}
设 $X$ 是完美域 $\Bbbk$ 上的有限型仿射概形, $G$ 是作用在 $X$ 上的约化 $\Bbbk$-群. 本习题将证明, 对每个点 $x\in X(\Bbbk)$, 存在定义在 $\Bbbk$ 上的单参数子群 $\lambda:\mathbb G_m\to G$, 使 $\lim_{t\to0}\lambda(t)\cdot x$ 具有闭 $G$-轨道. 另见 \cite[\S4]{Kem78}.
\begin{enumerate}[(1)]
\item 证明几何 Galois 群 $\operatorname{Gal}:=\operatorname{Gal}(\overline{\Bbbk}/\Bbbk)$ 作用在单参数子群集合 $\mathbb X_*(G_{\overline{\Bbbk}})$ 上, 且 $\mathbb X_*(G)=\mathbb X_*(G_{\overline{\Bbbk}})^{\operatorname{Gal}}$.
\item 证明 $\mathbb X_*(G_{\overline{\Bbbk}})$ 上存在 $\operatorname{Gal}$-不变长度 $\|{-}\|$.
\item 证明子集
\[
\begin{aligned}
&\{\lambda\in\mathbb X_*(G_{\overline{\Bbbk}})\mid\lim_{t\to0}\lambda(t)\cdot x\in X(\overline{\Bbbk})\text{ 存在}\},\\
&\{\lambda\in\mathbb X_*(G_{\overline{\Bbbk}})\mid\lim_{t\to0}\lambda(t)\cdot x\in G_{\overline{\Bbbk}}x_0\}
\end{aligned}
\]
是 $\operatorname{Gal}$-不变的, 其中 $G_{\overline{\Bbbk}}x_0$ 是 $\overline{G_{\overline{\Bbbk}}x}$ 中唯一的闭轨道. 此外, 证明若 $V$ 和 $W$ 是如~\bookeqref{8.6.8} 中的 $G$-表示, 则函数 $m(i(x),\lambda)$ 和 $m(f(x),\lambda)$ 都是 $\operatorname{Gal}$-不变的.
\item 将定理~\bookref{8.6.5} 和定理~\bookref{8.6.6} 推广到 $\Bbbk$ 为完美域且 $x\in X(\Bbbk)$ 的情形.
\end{enumerate}
特别地, 若 $G$ 没有定义在 $\Bbbk$ 上的非平凡单参数子群, 则任意 $\Bbbk$-点的 $G$-轨道都是闭的.
\end{exercise}
最后记录 Kempf 最优不稳定化定理 (\bookref{8.6.5}) 的证明的如下推论. 它将在 HKKN 分层 (\bookref{8.7.1}) 的证明中起关键作用.
\begin{proposition}\label{src:8.6.16}
设 $G$ 是代数闭域 $\Bbbk$ 上的约化群, 在 $\mathbb X_*(G)$ 上给定长度 $\|{-}\|$. 设 $X=\Spec A$ 是带有 $G$ 作用的有限型仿射 $\Bbbk$-概形, 且具有唯一闭轨道 $Gx_0$. 固定极大环面 $T\subseteq G$. 存在有限多个单参数子群 $\lambda_1,\ldots,\lambda_n\in\mathbb X_*(T)$ 和数 $M_1,\ldots,M_n\in\mathbb R_{<0}$, 使对每个点 $x\in X(\Bbbk)$, 存在唯一 $i=1,\ldots,n$, 满足 $\lambda_i$ 是某个 $gx$, $g\in G$, 的最优不稳定化单参数子群, 且 $M_i=\mu(x,\lambda_i)/\|\lambda_i\|$.
\end{proposition}
\begin{proof}
使用定理~\bookref{8.6.5} 证明中的记号. 对 $x\in X(\Bbbk)$, 最优不稳定化单参数子群的唯一抛物子群由当 $g$ 遍历 $G(\Bbbk)$ 时的子集
\begin{align*}
\operatorname{State}_{gTg^{-1}}(i(x))&\simeq\operatorname{State}_T(i(gx))\subseteq\mathbb X_*(T),\\
\operatorname{State}_{gTg^{-1}}(f(x))&\simeq\operatorname{State}_T(f(gx))\subseteq\mathbb X_*(T)
\end{align*}
确定. 这些子集包含在由满足 $V_\chi\ne0$ 或 $W_\chi\ne0$ 的特征 $\chi\in\mathbb X_*(T)$ 组成的有限子集中. 因而 $T$ 中的最优不稳定化子群仅有有限多种可能.
\end{proof}
